\documentclass[10pt,a4paper]{amsart}
\usepackage[margin=19mm]{geometry}
\usepackage{amsmath,amssymb,mathtools}
\usepackage[scr=rsfs,cal=euler]{mathalfa}
\usepackage{xfrac}
\usepackage[unicode,hidelinks,bookmarks=false]{hyperref}
\newcommand{\ie}{i.e.\ }
\newcommand{\cf}{cf.\ }
\newcommand{\resp}{respectively\ }
\newcommand{\Subsection}{\subsection}
\providecommand{\nobreakdash}{\nobreak\hbox{-}}
\setlength{\emergencystretch}{2em}
\multlinegap0pt
\usepackage{bm}
\usepackage{bbm}
\usepackage{dsfont}
\usepackage{xspace}
\usepackage{cancel}
\usepackage{mathtools}
\usepackage{amsthm}
\makeatletter
\@ifundefined{corr}{\newtheorem{corr}{Corollary}}{}
\@ifundefined{lemma}{\newtheorem{lemma}{Lemma}}{}
\@ifundefined{remark}{\newtheorem*{remark}{Remark}}{}
\makeatother
\title[Regularity Conditions for Critical Point Convergence]{Regularity Conditions for Critical Point Convergence}
\author{Thomas J. Maullin-Sapey, Samuel Davenport}
\date{Source: arXiv:2507.01854v2 (CC BY 4.0)}
\begin{document}
\maketitle

\noindent\emph{Source and licence.} Regularity Conditions for Critical Point Convergence. Version arXiv:2507.01854v2 (CC BY 4.0), distributed under the Creative Commons Attribution 4.0 International licence, as recorded in the arXiv licence metadata for this exact version and shown on the abstract page https://arxiv.org/abs/2507.01854v2. Full rights notice: licenses/arxiv-2507.01854-CC-BY-4.0.txt.\par\smallskip

\noindent\textit{Editorial scope.} Counted unit: the Remark whose source label is \texttt{None} in the article (source0.tex lines 1495--1501, source label \texttt{None}), delivered together with the article's own complete original proof of it (source0.tex lines 1502--1522, one \texttt{proof} environment of 21 lines). No mathematics is written, repaired or re-derived: statement and proof are the article's own text line for line. Required context retained uncounted: lem:lipschitz (lines 1053--1055); lem:morse (lines 1230--1234); corr:morse\_constants (lines 1486--1492). Every definition, display and label of those lines is retained unchanged. Omitted from this package and not redistributed: 1--1052, 1056--1229, 1235--1485, 1493--1494, 1523--1533. Nothing inside a retained excerpt is omitted or altered: every retained body is delivered exactly as the article's lines are written, so no omission receipt of any kind accompanies this package. Citations: the retained excerpts carry no citation command at all, so no bibliography entry of the article is transcribed and no reference is redistributed. Reference adaptation: none was needed and none was made; every reference inside the delivered unit resolves inside it, so no unresolved reference or citation is produced. Printed slips kept literal: no printed word, symbol, hypothesis or number of the counted statement or of its proof was altered. Layout and class: the article's own class is \texttt{article} with packages including amsmath, amssymb, amsthm, bm, mathtools, appendix, dsfont, floatflt, verbatim, graphicx, rotating, xcolor, caption, natbib; the wrapper is the lane's pinned \texttt{amsart} editorial wrapper at 10pt on A4, which restates the theorem-like environments the retained text needs and adds the article's own macro definitions that the retained text needs (none), with extra wrapper packages (bm, bbm, dsfont, xspace, cancel, mathtools). The delivered numbering is the unit's own. Assets: no retained span contains a figure environment, a graphics-inclusion command, a PDF-page inclusion, a table or a TikZ picture; no image file is copied into this package, the package declares no asset dependency and no image file is redistributed with it. Licence: version arXiv:2507.01854v2 of this article is distributed under the Creative Commons Attribution 4.0 International licence, as recorded in the arXiv licence metadata for this exact version and shown on the abstract page https://arxiv.org/abs/2507.01854v2. A full rights notice is delivered at licenses/arxiv-2507.01854-CC-BY-4.0.txt. Characters: every retained excerpt is pure ASCII, so the delivered engine, XeLaTeX with its default Latin Modern text font, reports no missing character.
    \begin{lemma}\label{lem:lipschitz}
        Suppose $\phi:S\rightarrow \mathbb{R}$ is $C^2$ with $H_\phi(0)=0$. Then, for all $L>0$, there exists an $R_0>0$ satisfying $L\geq \sup_{x\in B_{R_0}}||H_\phi(x)||$, such that $\nabla \phi$ is Lipschitz on $B_{R_0}$ with constant $L$.
    \end{lemma}

\begin{lemma}[The Morse Lemma]\label{lem:morse} Suppose that $S$ is a $C^{2}$ manifold and let $f:S\rightarrow \mathbb{R}$ be Morse. Assume $f$ has a Morse point at the origin and, as in the main text, denote $H:=H_f(0)$. Then, there exists $r >0$ depending only upon $||H||$ and $||H^{-1}||$ and a zero-preserving $C^1$ bi-Lipschitz homeomorphism $\Gamma(x)$ defined on the ball $B_r$ such that:
\begin{equation}\nonumber
   f(\Gamma(x)) = \frac{1}{2}x'Hx.
\end{equation}
\end{lemma}

    \begin{corr}[Constants of the Morse Lemma]\label{corr:morse_constants}
        Suppose $r$ is a positive constant satisfying:
        \begin{equation}\nonumber
            \sup_{x\in B_r}||H_f(x)-H|| < ||H^{-1}||^{-1}\bigg[(1-m) \land  \frac{m^2\ln(2)}{6||H||||H^{-1}||}\bigg] \quad \text{ and }\quad r < \frac{m}{4||H||||H^{-1}||}
        \end{equation}
        for some $m\in (0,1)$, then $r$ satisfies the Morse Lemma. That is, $\Gamma$ is a well-defined bi-Lipschitz diffeomorphism on $B_{r}$.
    \end{corr}

\begin{remark}
    It is worth emphasizing that the conclusion of Corollary \ref{corr:morse_constants} holds \text{for any} $m\in (0,1)$. For instance, choosing $m=\sfrac{1}{2}$, we see that the conditions reduce to:
        \begin{equation}\nonumber
            \sup_{x\in B_r}||H_f(x)-H|| < \frac{1}{2||H^{-1}||} \land  \frac{\ln(2)}{24||H||||H^{-1}||^2} \quad \text{ and }\quad r < \frac{1}{8||H||||H^{-1}||}.
        \end{equation}
    It follows that the value of $r$ depends only on the constants $||H||$ and $||H^{-1}||$, and the behavior of $H_f$ near the origin.
\end{remark}

\begin{proof}
    By tracking the constants in the proof of Lemma \ref{lem:morse}, we see that $\Gamma_t$ is a well-defined bi-Lipschitz diffeomorphism if the following criteria are met:
    \begin{equation}\nonumber
       \textbf{(1)}~~ L \leq ||H^{-1}||^{-1}, \quad \textbf{(2)}~~ 2-e^{a_1} > 0, \quad \text{ and } \quad \textbf{(3)}~~ r \leq \frac{c}{C}e^{-a_1}.
    \end{equation}
    where $L$ is the Lipschitz constant of $\phi$ on $B_r$, and $a_1, c$ and $C$ are as defined previous.\\
    \\
    Let $m\in(0,1)$ be arbitrary and let $r$ satisfy the statement of the corollary. For notational ease, let $\epsilon=m||H^{-1}||^{-1}$. By Lemma \ref{lem:lipschitz} and the assumption of the corollary, $L\leq\sup_{x\in B_r}||H_f(x)-H||<||H^{-1}||^{-1}-\epsilon$ and thus, the first criterion holds. Further, we also have that:
    \begin{equation}\nonumber
        L < \frac{\epsilon^2\text{ln}(2)}{6||H||} < \frac{\epsilon^2\text{ln}(2)}{2||H^{-1}||^{-1} + 3||H|| + L} = \frac{\epsilon^2\text{ln}(2)}{2(||H^{-1}||^{-1}-L) + 3(||H|| + L)},
    \end{equation}
    as, by common matrix identities and the above $L<||H^{-1}||^{-1}\leq||H||$. Thus, as $\epsilon < ||H^{-1}||^{-1}-L=c$, we have that:
    \begin{equation}\nonumber
        L < \frac{\epsilon^2\text{ln}(2)}{2c + 3C} < \frac{c^2\text{ln}(2)}{2c + 3C}.
    \end{equation}
    Rearranging and noting the definition of $a_1$, we have that $a_1<\text{ln}(2)$ and thus $2-e^{a_1}>0$, i.e. the second criterion holds.  Finally, through similar logic to the above, we have that:
    \begin{equation}\nonumber
        r < \frac{m||H^{-1}||^{-1}}{4||H||} \leq \frac{||H^{-1}||^{-1}-L}{2(||H||+L)} = \frac{1}{2}\cdot\frac{c}{C} \leq \frac{c}{C} e^{a_1},
    \end{equation}
    where the final inequality follows from the previous. Thus condition (3) holds. The result now follows.
\end{proof}

\end{document}
